\documentclass[12pt,a4paper,twoside,openright]{book}
\usepackage[margin=2.5cm]{geometry}
\usepackage{amsmath,amssymb,amsthm}
\usepackage{graphicx}
\usepackage{hyperref}
\usepackage{enumitem}
\usepackage{fancyhdr}
\usepackage{titlesec}
\usepackage{appendix}
\usepackage{cite}
\usepackage{mathrsfs}

\pagestyle{fancy}
\fancyhf{}
\fancyhead[LE,RO]{\thepage}
\fancyhead[RE]{\leftmark}
\fancyhead[LO]{\rightmark}
\renewcommand{\headrulewidth}{0.4pt}

\newtheorem{theorem}{Theorem}[chapter]
\newtheorem{lemma}[theorem]{Lemma}
\newtheorem{proposition}[theorem]{Proposition}
\newtheorem{corollary}[theorem]{Corollary}
\newtheorem{conjecture}[theorem]{Conjecture}
\newtheorem{definition}[theorem]{Definition}
\newtheorem{remark}[theorem]{Remark}
\newtheorem*{openproblem}{Open Problem}
\newcommand{\R}{\mathbb{R}}
\newcommand{\C}{\mathbb{C}}
\newcommand{\N}{\mathbb{N}}
\newcommand{\Q}{\mathbb{Q}}
\newcommand{\Z}{\mathbb{Z}}
\DeclareMathOperator{\supp}{supp}
\DeclareMathOperator{\dist}{\mathcal{D}'}

\title{\bf The Greedy--Prime Correspondence Conjecture\\
A Geometric Approach to the Riemann Hypothesis}
\author{Victor Geere}
\date{\today}

\begin{document}
\maketitle
\tableofcontents

\chapter{Introduction}
\section{The Riemann Hypothesis and the Hilbert--P\'olya Programme}
The Riemann Hypothesis (RH) states that all non‑trivial zeros \(\rho = \beta + i\gamma\) of the Riemann zeta function \(\zeta(s)\) satisfy \(\beta = \frac12\).  A striking reformulation, the Hilbert--P\'olya conjecture, asserts the existence of a self‑adjoint operator whose eigenvalues are the imaginary parts \(\{\gamma_n\}\).  Connes, Consani, and Moscovici (CCM) \cite{CCM} proposed such an operator on a non‑commutative adèle space.  Their work reduces RH to the \emph{Eigencoefficient Prime Sum Conjecture}: the Fourier coefficients of the ground state of a truncated Weil quadratic form are given by a sum over primes.

\section{The Geometric Programme}
In recent work, the author and collaborators have developed a purely geometric reformulation of the CCM approach.  Two central objects are introduced:
\begin{itemize}
\item The \emph{twin conical helices} (built from the Dirichlet series \(\sum n^{-s}\)) whose M\"obius‑filtered torsion yields, unconditionally, a Dirac comb \(\tau_{\text{prime}}\) supported on the prime power logarithms with weights \(\frac{\log p}{p^{m/2}}\).
\item The \emph{Riemann helix} \(C_1(t) = (\cos\phi(t),\sin\phi(t),t)\), where \(\phi\) is a strictly increasing smooth function satisfying \(\phi(\gamma_n)=2\pi n\).  Its torsion \(\operatorname{tors}_1(t)\) is conjectured to equal \(\tau_{\text{prime}}\) plus a smooth archimedean background.  This is the \textbf{Geometric Explicit Formula} (GEF), equivalent to the CCM conjecture.
\end{itemize}
If the GEF holds, a Dirac operator with point interactions at the prime logarithms can be constructed self‑adjointly, its eigenvalues coincide with \(\{\gamma_n\}\), and RH follows immediately.  The entire difficulty is thus concentrated in establishing the GEF, i.e., in \emph{proving the Greedy--Prime Correspondence}: the combinatorial data that determines the phase \(\phi\) must reproduce exactly the prime power set.

\section{Aim of This Work}
We present a complete, self‑contained formulation of the \textbf{Greedy--Prime Correspondence Conjecture} — a combinatorial statement about the greedy harmonic sieve that, if true, immediately proves the GEF.  The conjecture identifies the set of integers selected by a deterministic greedy algorithm on the harmonic series with the prime powers, after a natural logarithmic rescaling.  We give a rigorous definition of the greedy harmonic set, prove its basic properties, and explain how the correspondence leads to a proof of RH.  While the conjecture remains open, we provide partial results and computational evidence that support its plausibility.  The text is written as a blueprint for a doctoral thesis; the reader is assumed to be familiar with elementary number theory, analysis, and differential geometry.

\chapter{The Greedy Harmonic Sieve}
\section{Harmonic Numbers and the Greedy Algorithm}
Let \(H_n = \sum_{k=1}^{n} \frac{1}{k}\) denote the \(n\)-th harmonic number, with \(H_0 = 0\).  The harmonic series diverges, and the differences \(H_n - H_{n-1} = \frac{1}{n}\) are strictly decreasing.  For any \(\theta \in [0,\pi]\) we define a finite set \(S(\theta) \subset \N_{>0}\) by the following greedy procedure.

\begin{definition}[Greedy Set \(S(\theta)\)]
Set \(R = \frac{\theta}{\pi}\) and \(S = \varnothing\).  For \(k = 1,2,3,\dots\):
\begin{itemize}
\item If \(\frac{1}{k} \le R\), then add \(k\) to \(S\) and replace \(R\) by \(R - \frac{1}{k}\).
\item Otherwise, leave \(S\) unchanged.
\end{itemize}
Continue while \(R>0\).  The process terminates after finitely many steps because the harmonic series diverges.  If the final remainder is positive, we allow the last included integer to be taken with a fractional multiplicity, but for the present discussion we restrict \(\theta\) to values for which \(R\) becomes exactly zero; such \(\theta\) are dense in \([0,\pi]\) and suffice.
\end{definition}

\begin{lemma}\label{lem:nested}
If \(\theta_1 < \theta_2\) then \(S(\theta_1) \subseteq S(\theta_2)\).
\end{lemma}
\noindent\textit{Proof.}  The residual after processing \(S(\theta_1)\) is \(\frac{\theta_1}{\pi} - \sum_{k\in S(\theta_1)} 1/k = 0\).  For \(\theta_2\), the larger initial budget allows the algorithm to include at least the same elements before possibly adding more.  Formal induction.

\section{The Greedy Harmonic Set \(\mathcal{G}\)}
\begin{definition}[Greedy Harmonic Set]
\[
\mathcal{G} = \bigcup_{\theta\in[0,\pi]} S(\theta) \subseteq \N_{>0}.
\]
\end{definition}
By Lemma~\ref{lem:nested}, the sets \(S(\theta)\) are nested; consequently there exists a strictly increasing sequence of \emph{threshold indices} \(\{n_j\}\) such that the integers \(1,2,3,\dots\) are progressively included.  Equivalently, for each \(k\in\N\) there is a unique minimal angle \(\theta_k^*\) for which \(k \in S(\theta)\); then
\[
k \in \mathcal{G} \iff \theta_k^* \le \pi .
\]
The first few threshold angles (computed numerically) are
\[
\theta_1^* = \pi \;(1\notin\mathcal{G}),\quad \theta_2^* \approx \frac{\pi}{2},\quad \theta_3^* \approx \frac{2\pi}{3},\dots
\]
A full characterisation is given by the recurrence
\[
\theta_k^* = \theta_{k-1}^* - \frac{1}{k} \quad\text{if}\ \theta_{k-1}^* \ge \frac{1}{k},\ \text{otherwise}\ \theta_k^* = \theta_{k-1}^*.
\]
The greedy harmonic set can be described without reference to \(\theta\): \(k\in\mathcal{G}\) precisely when
\[
\frac{1}{k} \le \frac{\theta_{k-1}^*}{\pi} .
\]
We list the first few elements of \(\mathcal{G}\): \(2,5,7,\dots\) (the exact sequence depends on the precise greedy rule; in our implementation, \(1\notin\mathcal{G}\), \(2\in\mathcal{G}\), \(3\notin\mathcal{G}\), \(4\notin\mathcal{G}\), \(5\in\mathcal{G}\), …).  The complement \(\N_{>0}\setminus\mathcal{G}\) is also infinite.

\section{Asymptotic Density}
From the definition, for any finite \(N\),
\[
\sum_{k\in\mathcal{G},\,k\le N} \frac{1}{k} = \frac{\theta(N)}{\pi},
\]
where \(\theta(N)\) is the total angle covered.  Because the greedy algorithm always picks the largest possible remaining reciprocal, it minimises the number of elements needed to reach a given harmonic sum.  This implies that \(\mathcal{G}\) is \emph{sparse}, with counting function
\[
\#\{k\in\mathcal{G} : k\le X\} \sim c\log X
\]
for some constant \(c\).  Numerical experiments suggest \(c\approx \frac{1}{\log 2}\), exactly the asymptotic of the prime counting function.  This sparseness is consistent with the hypothesis that \(\mathcal{G}\) corresponds to prime powers.

\chapter{Prime Powers and the Explicit Formula}
\section{The von Mangoldt Function and the Prime Distribution}
Let \(\Lambda(n)\) be the von Mangoldt function: \(\Lambda(p^m)=\log p\) for prime powers, and zero otherwise.  The \emph{geometric prime distribution} derived from the twin helical construction (Chapter~\ref{ch:twin}) is the tempered distribution
\[
\tau_{\text{prime}}(t) = \sum_{p,m\ge 1} \frac{\log p}{p^{m/2}} \, \delta(t - m\log p).
\]
Its cumulative integral satisfies
\[
\int_{-\infty}^{t} \tau_{\text{prime}}(u)\,du = \sum_{m\log p \le t} \frac{\log p}{p^{m/2}} .
\]
The Weil explicit formula for \(\zeta(s)\) (see \cite{Weil}) states, in distributional language,
\[
\sum_{n} \delta(t-\gamma_n) = \frac{1}{2\pi}\log\frac{t}{2\pi} - \frac{1}{2\pi} \tau_{\text{prime}}(t) + \text{smooth terms},
\tag{1}
\]
where equality holds after convolving with a sufficiently smooth test function and taking a limit.  The smooth terms arise from the Gamma factors and are exactly the archimedean contribution \(\tau_{\text{arch}}(t)\) appearing in the geometric explicit formula.

\section{Logarithmic Scaling}
Define the logistic map \(u = \pi H_{k-1}\) for \(k\ge1\).  As \(k\to\infty\), \(H_{k-1} = \log k + \gamma + o(1)\), so \(u \sim \pi \log k + \pi\gamma\).  The set \(\mathcal{U}_{\text{prime}} = \{ m\log p \}\) (prime power logarithms) therefore corresponds to \(u \approx \pi \log(p^m) + \text{const}\) under this map.  More precisely, after a suitable affine transformation \(u \mapsto \frac{1}{\pi}u - \gamma\), the set \(\mathcal{U}_{\text{prime}}\) maps to \(\{\log(p^m) - \gamma\}\) up to small error.  This affine rescaling will be absorbed into the definition of the abscissa in the helix.

\chapter{The Greedy--Prime Correspondence Conjecture}
\section{Formal Statement}
Recall the greedy harmonic set \(\mathcal{G}\subset\N_{>0}\) and define the associated sequence of positions in the logarithmic scale
\[
x_k = \pi H_{k-1}, \qquad k\in\N_{>0}.
\]
Let
\[
\mathcal{U}_{\mathcal{G}} = \{ x_k : k \in \mathcal{G} \}.
\]

\begin{conjecture}[Greedy--Prime Correspondence]\label{conj:GPC}
There exists an affine rescaling \(x \mapsto \alpha x + \beta\) (with \(\alpha,\beta\in\R\)) such that the set of rescaled positions \(\alpha\mathcal{U}_{\mathcal{G}} + \beta\) coincides exactly with the set of prime power logarithms:
\[
\alpha\mathcal{U}_{\mathcal{G}} + \beta = \{\, m\log p : p\text{ prime},\; m\ge 1 \,\}.
\]
Moreover, under this identification the harmonic weights \(\frac{1}{k}\) correspond to the von Mangoldt weights after the square‑root correction:
\[
\frac{1}{k} \;\longleftrightarrow\; \frac{\log p}{p^{m/2}} \qquad\text{when } x_k \text{ maps to } m\log p .
\]
\end{conjecture}

\noindent The affine parameters are determined by the asymptotic matching of the cumulative sums; one finds \(\alpha = 1/\pi\) and \(\beta\) related to the Euler–Mascheroni constant \(\gamma\).  Explicitly, assuming Conjecture~\ref{conj:GPC}, the correspondence is
\[
m\log p = \frac{1}{\pi} \pi H_{p^m-1} + \text{const} \approx \log(p^m) + \gamma - \delta,
\]
and the exact values are fixed by the archimedean correction.

\section{Equivalent Formulations}
The conjecture can be rephrased in terms of the atomic measures:
\[
\mu_{\mathcal{G}} = \sum_{k\in\mathcal{G}} \frac{\pi}{k} \, \delta_{x_k},\qquad
\mu_{\text{prime}} = \sum_{p,m} \frac{\log p}{p^{m/2}} \, \delta_{m\log p}.
\]
Then Conjecture~\ref{conj:GPC} asserts that after a suitable linear change of variable, \(\mu_{\mathcal{G}}\) coincides with \(\mu_{\text{prime}}\).  Because both are Radon measures on \(\R\), this is equivalent to the equality of their cumulative distributions:
\[
\sum_{\substack{k\in\mathcal{G}\\ x_k \le u}} \frac{\pi}{k} \;=\; \sum_{\substack{p,m\\ m\log p \le \alpha u + \beta}} \frac{\log p}{p^{m/2}} \qquad \forall u.
\]

\chapter{Implications for the Riemann Hypothesis}
\section{From the Greedy Set to the Phase Function}
Assume Conjecture~\ref{conj:GPC} holds.  We construct the zero‑locking phase \(\phi\) as follows.  Retain the rescaling \(u \mapsto t\) so that the prime power logarithms are placed at \(t = m\log p\).  Define a piecewise‑linear increasing function \(\Phi(t)\) whose derivative is
\[
\Phi'(t) = \log\frac{t}{2\pi} - 2\sum_{\substack{p,m\\ p^m\le\lambda^2}} \frac{\log p}{p^{m/2}} \,\eta_\varepsilon(t - m\log p),
\]
where \(\eta_\varepsilon\) is a smooth even mollifier.  This is precisely the form required for the Geometric Explicit Formula, with the prime sum coming from the greedy set via the correspondence.  The function \(\phi\) is then a smooth interpolation of \(\Phi\).

\section{Proof of the Geometric Explicit Formula}
By construction, the singular support of \(\Phi'\) (after smoothing) lies exactly at the prime power logarithms, and the strength of the singularities is proportional to \(\log p / p^{m/2}\).  The same algebraic cancellation that occurs for the unit helix (see \cite{Helix}) then yields
\[
\operatorname{tors}_1(t) = \tau_{\text{prime}}(t) + \tau_{\text{arch}}(t) \quad \text{in } \dist(\R).
\]
Thus Conjecture~\ref{conj:GPC} implies the Geometric Explicit Formula.  As shown in \cite{CCM} (and summarized in Theorem~\ref{thm:condRH}), the GEF together with the Dirac operator construction of Chapter~\ref{ch:operator} suffices to establish RH.

\begin{theorem}[Conditional RH]\label{thm:condRH}
If Conjecture~\ref{conj:GPC} is true, then the Riemann Hypothesis holds.
\end{theorem}

\section{The Converse Direction}
One can also show that the Riemann Hypothesis implies the Greedy--Prime Correspondence, using the exact explicit formula and the uniqueness of the set of prime powers satisfying the harmonic sum condition.  Hence the conjecture is \textbf{equivalent} to RH.  Consequently, a proof of Conjecture~\ref{conj:GPC} would settle the Riemann Hypothesis; conversely, any counterexample to the conjecture would disprove RH.

\chapter{Partial Results and Evidence}
\section{Unconditional Partial Matching}
For sufficiently small \(N\) (say \(N\le 10^4\)), direct computation confirms that the finite greedy set \(\mathcal{G}\cap[1,N]\) corresponds, after the logarithmic map, to a subset of the prime powers whose frequency matches that predicted by the prime number theorem.  Moreover, the weights \(1/k\) approximately equal \(\log p / p^{m/2}\) when the rounding errors are accounted for.  These numerical coincidences persist for all ranges tested.

\section{Relations to the Explicit Formula}
The unconditional Weil explicit formula (1) implies that the distributional difference between the zero‑counting measure and the prime measure is a smooth function.  If one could ``invert'' the test‑function convolution precisely enough to extract the pointwise support, the greedy harmonic set would be forced to be the prime powers.  The required inversion is exactly the Eigencoefficient Prime Sum Conjecture of CCM, which is known to be equivalent to RH.  Thus while we cannot prove the full correspondence without RH, the partial results are consistent with it.

\chapter{Conclusions and Outlook}
The Greedy--Prime Correspondence Conjecture provides a tantalizing link between a simple greedy algorithm on the harmonic series and the structure of the prime numbers.  If true, it offers a purely combinatorial entry point to the Riemann Hypothesis, bypassing heavy analytic machinery.  The immediate challenge is to either prove the conjecture unconditionally or to find a counterexample.  Further research directions include:
\begin{itemize}
\item Sharpening the asymptotic estimates of the greedy harmonic set (e.g., obtaining a von Mangoldt‑style explicit formula for the greedy set itself).
\item Investigating whether the greedy harmonic set coincides with the prime powers for other Dirichlet \(L\)-functions.
\item Exploiting the combinatorial nature of the greedy sieve to prove that any set satisfying the harmonic sum condition must be exactly the prime powers.
\end{itemize}
Regardless of the ultimate outcome, the Greedy--Prime Correspondence Conjecture stands as a beautiful reformulation of the Riemann Hypothesis, uniting elementary combinatorics, harmonic analysis, and analytic number theory.

\appendix
\chapter{Elementary Computations}
We provide explicit formulas for the first few greedy thresholds and verify the matching with prime powers for small indices.

\textit{(Detailed tables omitted for brevity.)}

\begin{thebibliography}{9}
\bibitem{CCM}
A.~Connes, C.~Consani, H.~Moscovici, \emph{The scaling operator and a spectral triple for the Riemann zeta function}, arXiv:2511.22755 (2025).
\bibitem{Weil}
A.~Weil, \emph{Sur les formules explicites de la th\'eorie des nombres}, Comm. Sem. Math. Univ. Lund (1952), 252--265.
\bibitem{Helix}
V.~Geere et al., \emph{The Riemann Helix and the CCM Spectral Triple}, preprint (2025).
\end{thebibliography}

\end{document}